\chapter*{Sources, Grades, and How to Read Further}
\addcontentsline{toc}{chapter}{Sources, Grades, and How to Read Further}

This condensed edition carries no bibliography of its own, deliberately. Every figure, quotation, and forecast in it is taken from the full edition of the same version and date, which cites each to its source---primary wherever possible---and grades it: \textbf{[V]} independently verified, \textbf{[P]} primary-source claim, \textbf{[A]} analyst estimate, \textbf{[M]} author-modeled result, \textbf{[R]} roadmap, \textbf{[S]} scenario judgment. Readers who wish to check a number should go to the corresponding chapter of the full edition, marked in this text as \full{n}, and from there to its evidence ledger, which is the authoritative record and the fastest route to the argument's weakest link.

The full edition also states its own known weaknesses, and a reader of the abridgement should know them too: many mid-2026 details rest on secondary aggregators; the attribution of the LX2 processor's manufacturing is analyst inference; the capacity-first memory design is a modeled architecture with no measured realization; the model cards remain incomplete in the fields that matter most for procurement, tokens and wall-clock per solved task; the national sizing is measurement-anchored for one country only; and the architectural thesis that the CPU corner's breadth compounds faster than the accelerator corner's depth is a judgment, not a measurement, and it is the judgment most likely to be shown wrong. Where this edition says ``the book argues,'' it means exactly that.

The reproducible models behind the numbers---the appliance-versus-rack cost model and the AI-for-science value funnel---ship with the source of the full edition so that a reader can change an input and re-derive the conclusion; that, rather than any single number, is the point of publishing them.

Both editions are living documents. The title-page date is the currency date; the claim register and dashboard are re-graded at every revision; corrections are welcome and will be credited, and a dated primary source moves the argument faster than an argument does.
